% GENERATED FILE. Edit canonical lessons, then run npm run generate:books.
% canonical-source-hash: fnv1a64:3f712c920b9484d9
% canonical-lessons: KA-C242-oggarane-haku, KA-C242-nenesu, KA-C242-hindu, KA-C242-addu, KA-C242-hari

\chapter{To temper, To soak, To wring, To dip, To tear}
\label{ch:ka-to-temper-to-soak-to-wring-to-dip-to-tear}
\hlchaptermodality{242}

\begin{chapteropening}
\textbf{By the end of this chapter:} \emph{I can say to temper, to soak, to wring, to dip and to tear.}

Complete the last lesson of chapter 242: I can say to temper, to soak, to wring, to dip and to tear.
\end{chapteropening}

\section[\texorpdfstring{\kn{ಒಗ್ಗರಣೆ} \kn{ಹಾಕು} (oggaraṇe hāku)}{ಒಗ್ಗರಣೆ ಹಾಕು (oggaraṇe hāku)}]{\kn{ಒಗ್ಗರಣೆ} \kn{ಹಾಕು} (oggaraṇe hāku) — to temper (fry spices in hot oil)}

\label{lesson:KA-C242-oggarane-haku}



\begin{quote}
Before the new one: say the Kannada for to find out, then the Kannada for to borrow.
\end{quote}

\subsection*{You'll want to know: \kn{ಒಗ್ಗರಣೆ} \kn{ಹಾಕು}}
\textbf{\kn{ಒಗ್ಗರಣೆ} \kn{ಹಾಕು}} — \emph{oggaraṇe hāku} — \textquotedblleft{}to temper (fry spices in hot oil)\textquotedblright{}.

\begin{grammarlens}[title={What you've built}]
One more word for this chapter.
\end{grammarlens}

\subsection*{Your turn}
\begin{itemize}
\raggedright
  \item \emph{Say it:} \emph{oggaraṇe hāku}
  \item \emph{Say it:} \emph{oggaraṇe hāku}, once more
  \item \emph{Say it:} say \emph{oggaraṇe hāku}
  \item \emph{Recall:} say the Kannada for to find out, then the Kannada for to borrow, then say \emph{oggaraṇe hāku} again
\end{itemize}

\begin{tcolorbox}[breakable,colback=teal!4,colframe=teal!35!black,title={Before you move on}]
What does \kn{ಒಗ್ಗರಣೆ} \kn{ಹಾಕು} mean? (\textquotedblleft{}To temper (fry spices in hot oil)\textquotedblright{}.) Say it once more.
\end{tcolorbox}
\section[\texorpdfstring{\kn{ನೆನೆಸು} (nenesu)}{ನೆನೆಸು (nenesu)}]{\kn{ನೆನೆಸು} (nenesu) — to soak}

\label{lesson:KA-C242-nenesu}



\begin{quote}
Before the new one: say the Kannada for to borrow, then the Kannada for to temper.
\end{quote}

\subsection*{You'll want to know: \kn{ನೆನೆಸು}}
\textbf{\kn{ನೆನೆಸು}} — \emph{nenesu} — \textquotedblleft{}to soak\textquotedblright{}.

\begin{grammarlens}[title={What you've built}]
One more word for this chapter.
\end{grammarlens}

\subsection*{Your turn}
\begin{itemize}
\raggedright
  \item \emph{Say it:} \emph{nenesu}
  \item \emph{Say it:} \emph{nenesu}, once more
  \item \emph{Say it:} say \emph{nenesu}
  \item \emph{Recall:} say the Kannada for to borrow, then the Kannada for to temper, then say \emph{nenesu} again
\end{itemize}

\begin{tcolorbox}[breakable,colback=teal!4,colframe=teal!35!black,title={Before you move on}]
What does \kn{ನೆನೆಸು} mean? (\textquotedblleft{}To soak\textquotedblright{}.) Say it once more.
\end{tcolorbox}
\section[\texorpdfstring{\kn{ಹಿಂಡು} (hiṇḍu)}{ಹಿಂಡು (hiṇḍu)}]{\kn{ಹಿಂಡು} (hiṇḍu) — to wring, to squeeze out}

\label{lesson:KA-C242-hindu}



\begin{quote}
Before the new one: say the Kannada for to temper, then the Kannada for to soak.
\end{quote}

\subsection*{You'll want to know: \kn{ಹಿಂಡು}}
\textbf{\kn{ಹಿಂಡು}} — \emph{hiṇḍu} — \textquotedblleft{}to wring, to squeeze out\textquotedblright{}.

\begin{grammarlens}[title={What you've built}]
One more word for this chapter.
\end{grammarlens}

\subsection*{Your turn}
\begin{itemize}
\raggedright
  \item \emph{Say it:} \emph{hiṇḍu}
  \item \emph{Say it:} \emph{hiṇḍu}, once more
  \item \emph{Say it:} say \emph{hiṇḍu}
  \item \emph{Recall:} say the Kannada for to temper, then the Kannada for to soak, then say \emph{hiṇḍu} again
\end{itemize}

\begin{tcolorbox}[breakable,colback=teal!4,colframe=teal!35!black,title={Before you move on}]
What does \kn{ಹಿಂಡು} mean? (\textquotedblleft{}To wring, to squeeze out\textquotedblright{}.) Say it once more.
\end{tcolorbox}
\section[\texorpdfstring{\kn{ಅದ್ದು} (addu)}{ಅದ್ದು (addu)}]{\kn{ಅದ್ದು} (addu) — to dip}

\label{lesson:KA-C242-addu}



\begin{quote}
Before the new one: say the Kannada for to soak, then the Kannada for to wring.
\end{quote}

\subsection*{You'll want to know: \kn{ಅದ್ದು}}
\textbf{\kn{ಅದ್ದು}} — \emph{addu} — \textquotedblleft{}to dip\textquotedblright{}.

\begin{grammarlens}[title={What you've built}]
One more word for this chapter.
\end{grammarlens}

\subsection*{Your turn}
\begin{itemize}
\raggedright
  \item \emph{Say it:} \emph{addu}
  \item \emph{Say it:} \emph{addu}, once more
  \item \emph{Say it:} say \emph{addu}
  \item \emph{Recall:} say the Kannada for to soak, then the Kannada for to wring, then say \emph{addu} again
\end{itemize}

\begin{tcolorbox}[breakable,colback=teal!4,colframe=teal!35!black,title={Before you move on}]
What does \kn{ಅದ್ದು} mean? (\textquotedblleft{}To dip\textquotedblright{}.) Say it once more.
\end{tcolorbox}
\section[\texorpdfstring{\kn{ಹರಿ} (hari)}{ಹರಿ (hari)}]{\kn{ಹರಿ} (hari) — to tear, to rip}

\label{lesson:KA-C242-hari}



\begin{quote}
Before the new one: say the Kannada for to wring, then the Kannada for to dip.
\end{quote}

\subsection*{You'll want to know: \kn{ಹರಿ}}
\textbf{\kn{ಹರಿ}} — \emph{hari} — \textquotedblleft{}to tear, to rip\textquotedblright{}.

\begin{grammarlens}[title={What you've built}]
One more word for this chapter.
\end{grammarlens}

\subsection*{Your turn}
\begin{itemize}
\raggedright
  \item \emph{Say it:} \emph{hari}
  \item \emph{Say it:} \emph{hari}, once more
  \item \emph{Say it:} say \emph{hari}
  \item \emph{Recall:} say the Kannada for to wring, then the Kannada for to dip, then say \emph{hari} again
\end{itemize}

\begin{tcolorbox}[breakable,colback=teal!4,colframe=teal!35!black,title={Before you move on}]
What does \kn{ಹರಿ} mean? (\textquotedblleft{}To tear, to rip\textquotedblright{}.) Say it once more.
\end{tcolorbox}
